\chapter{基于任务分配的多传感器搜索与跟踪管控方法}

\section{引言}
第三章中提出的针对有限视域下多传感器管控方法通过构建加权搜索与跟踪目标函数，
实现了两类任务的联合优化。然而，该方法在很大程度上依赖于目标函数中的加权参数。
合适的加权参数才能使传感器需要在决策过程中同时兼顾搜索与跟踪，
而且这种方法只能得到加权目标函数下的最优动作，不能保证任一子任务在各自的目标函数下达到最优。
另外该方法对所有传感器的联合动作空间进行优化求解，
随着传感器规模的增加，决策空间呈指数增长，导致计算复杂度迅速上升，限制了算法的可扩展性。
针对上述问题，本文从任务结构重构的角度出发，对多传感器管控框架进行改进。
在决策优化之前，首先根据当前信息对传感器进行任务划分，使其分别承担搜索或跟踪任务，
从而避免单个传感器在两类任务之间进行权衡。通过这种方式，
不同传感器仅需针对各自任务进行策略优化，使得对应子任务的性能能够得到提升。
同时，优化过程由原来的全局联合搜索转化为在子任务内部进行决策求解，有效降低了计算复杂度。

在此基础上，本章仍针对有限视域下的多目标搜索与跟踪的多传感器管控问题，
构建多传感器管控方法。本章沿用第三章中的跟踪滤波、区域网格构建以及信息融合机制，
在此之上设计了任务分配策略，提出一种先任务分解再分组优化的两阶段管控方法。
通过仿真实验对所提方法进行验证，
结果表明该方法在保证搜索与跟踪性能的同时，
较上一章方法在计算效率与协同效果方面均有明显提升。

\subsection{两阶段多传感器管控框架}
在多目标优化问题中，不同目标之间往往存在冲突。对于单一优化目标而言，
任意两个候选解均可以通过统一指标进行比较，从而获得明确的最优解；然而，
在多目标搜索与跟踪问题中，针对未探测区域的搜索与对已知目标的精确跟踪需要共同
占用有限的传感器资源，这使得两类任务在优化过程中呈现出显著的竞争关系。
因此，将搜索性能与跟踪性能同时进行优化，本质上构成一个典型的多目标优化问题，
其最优解通常表现为不同目标之间的折衷结果，而难以使各子任务分别达到最优。

针对多目标优化问题，常见方法包括加权合成、约束转化及群智能优化等。
上一章通过构建加权目标函数，将搜索与跟踪任务统一到单目标函数下进行求解，
从而在一定程度上降低了问题复杂度。然而，该方法得到的解实质上属于帕累托有效解，
其性能依赖于权重参数的选取，难以保证搜索与跟踪两类任务分别达到理想最优状态。
同时，在该框架下，每个传感器在决策时均需同时考虑搜索与跟踪两类指标，
这在实际中并不总是合理：例如，当传感器视场内不存在已知目标时，
其参与跟踪指标优化的意义有限；
反之，对于已稳定跟踪目标的传感器，过度强调搜索收益也可能导致跟踪性能下降。

基于上述分析，本文从任务重构的角度出发，引入任务分配机制，
对搜索任务和跟踪任务进行分解，使所有节点在任何时刻都只独立执行一项任务。
通过该方式，将多目标联合优化问题，转化为独立的搜索子问题与跟踪子问题。
相较于对所有传感器进行联合优化，该方法既避免了不同任务间资源的竞争，
还降低了决策空间规模与计算复杂度，提升了多传感器系统的整体运行效率。

下面介绍基于任务分配的多传感器管控框架。

首先将传感器分为两类：S传感器为搜索任务传感器，任务目标是保持对已检测到目标的持续跟踪。
T传感器为跟踪任务传感器，任务目标是搜索未知区域以检测区域内可能存在的目标。
包含所有S传感器的集合记为$\mathcal{S}$，包含所有T传感器的集合记为$\mathcal{T}$
在任务分解阶段，首先在得到融合后的全局后验强度和区域网格图后，
根据融合结果为每一个传感器分配搜索或跟踪任务。
在分组优化阶段，首先根据3. 节的内容对传感器执行伪预测与伪更新，
之后针对执行搜索或跟踪任务的传感器集合进行动作决策。
基于任务分配的多传感器管控框架流程图如图4.1所示：
\begin{figure}[htbp]
    \centering
    \includegraphics[width=0.8\textwidth]{占位图.png}
    \caption{占位图}
    \label{fig:my_picture}
\end{figure}

\subsection{任务动态分配机制}
在多目标搜索与跟踪任务中，监视区域内存在目标的新生或消亡，目标可能随时出现在
某个传感器视域内，也可能随时消失。面对这种动态变化的环境，需要动态调整传感器
的任务。
任务的动态转换主要体现在两个方面：其一为响应目标出生与消亡的任务分配调整。
当目标离开传感器视域或消亡时，传感器应该及时从跟踪任务转为搜索任务。
反之，当新目标被发现时或传感器视域内有新生目标时，
传感器应该从搜索任务转为跟踪任务。其二。基于当前跟踪性能的任务分配调节。
若当前目标的跟踪精度未能达到预期阈值，
则需要增加执行跟踪任务的传感器数量，以提升跟踪精度，反之，若跟踪性能过剩时，
则应该分配更多的传感器执行搜索任务，以提升发现新目标的概率。
因此任务转换需要兼顾目标跟踪的精度需求与区域搜索的需求，
确保传感器网络能够灵活应对复杂多变的监视环境。

在初始时刻，传感器网络对监视区域内目标的数量、状态以及分布均未知。因此所有传感器节点均设为
S传感器。对于任意S节点，若连续$k$个时刻其视域内均有目标存在，将该节点从S传感器转为
T传感器。对于任意T节点，若连续$k$个时刻其视域内无目标存在，则将该节点从T传感器转为S传感器。
为了根据目标跟踪质量调整传感器任务类型，将所有被检测目标的平均跟踪误差作为目标跟踪质量的评判依据。
定义任务切换指标函数为：
\begin{equation}
\overline{\theta}^t=\sum_{i=1}^{N_k}\operatorname{tr}\left(WP_k^iW^\mathrm{T}\right)/N_k
\end{equation} 
其中，$W$为提取协方差矩阵中位置信息的矩阵，$N_k$为目标个数。

同时引入两个不确定度阈值$\theta_l$和$\theta_h$，
其中$\theta_l<\theta_h$。当平均不确定度$\overline{\theta}^t$超过较高的阈值$\theta_2$时，
表明当前的跟踪性能较差，需要增加更多的传感器参与跟踪任务，以提升跟踪精度。

跟踪性能最差的目标对应的高斯分量的索引为：
\begin{equation}
n=\max_{i}\mathrm{tr}\left(WP_{k}^{i}W^{\mathrm{T}}\right)
\end{equation} 
该目标的跟踪性能最差，需要优先为其分配更多的传感器资源以提升跟踪性能
。选择距离该目标最近的S传感器，转换为T传感器，执行跟踪任务。

当平均不确定度$\overline{\theta}^t$低于较低的阈值$\theta_1$时，表明当前的跟踪性能较好，
可以减少参与跟踪任务的传感器数量，将更多的传感器分配到搜索任务中，以提升新目标的发现能力。

跟踪性能最好的目标对应的高斯分量的索引为：
\begin{equation}
q=\min_{i}\mathrm{tr}\left(WP_{k}^{i}W^{\mathrm{T}}\right)
\end{equation} 
当跟踪该目标的传感器数量大于1时，对视域中包含该目标的T传感器进行评估，
选择距离该目标最远的一个T传感器，将其转换为S传感器，执行搜索任务。

当平均不确定度$\overline{\theta}^t$介于两个阈值之间时，表明当前的跟踪性能处于可接受范围内，
无需调整传感器的任务分配。

\subsection{针对跟踪任务的多传感器管控方法}
T传感器只执行对视域内目标的跟踪任务，对所有T传感器的后验强度信息融合后得到
融合后验强度，记为：

针对跟踪任务的目标函数表示如下：
\begin{equation}
J_{track}(\mathbf{c}_k)=\sum_{j=1}^{J}\mathrm{tr}\left(WP_{k+L}^j(\mathbf{c}_k)W^T\right)+\beta\cdot\left(\hat{N}_{k+L|k}-\hat{N}_{k+L}(\mathbf{c}_k)\right)
\end{equation}
其中，$\beta$是一个惩罚尺度因子,以防止有限视域传感器在协同决策时丢失对已发现目标的跟踪。$W$为提取目标位置分量的矩阵。
由于伪预测与伪更新过程基于PIMS生成，因此伪后验基数
$\hat{N}_{k+L}(\mathbf{c}_k)$一定低于先验预测基数$\hat{N}_{k+L|k}$。


\subsection{针对搜索任务的多传感器管控方法}
在搜索任务中，S传感器需要估计监视区域目标可能存在的概率，并根据需要搜索的区域生成有效搜索计划。
3.3节中已经对监视区域构建区域网格图，并通过预测与更新步骤估计监视区域目标可能存在的概率。
首先将第$i$个网格的目标存在概率记为该网格的搜索值$p_{sear}^{(i)}$。
针对处于搜索状态的传感器集合$\mathcal{S}$，构建整数线性规划模型，对目标存在概率高的区域进行全局
最优路径分配。
首先将任务区域建模为无向图$G=(\mathcal{V},\mathcal{E})$，
其中$\mathcal{V}$是区域网格节点集合，$\mathcal{E}$是节点之间的可达边集合。
以欧氏距离作为边$(i,l)$ 的移动代价 $c_{il}$。
将搜索值高于设定搜索阈值 的节点作为待搜索区域，记为$\bar{\mathcal{V}}$,
且$\bar{\mathcal{V}} \subseteq \mathcal{V}$。
为了计算每个传感器节点的最优路径，以最小化遍历成本，制定以下整数线性规划模型，
式（）为目标函数。
\begin{equation}
\min\sum_{j\in S}\sum_{(i,l)\in\mathcal{E}}c_{il}y_{il}^j
\end{equation} 
下面是该整数线性规划模型的约束条件：
确保未搜索区域$\bar{\mathcal{V}}$至少被一个传感器节点搜索：
\begin{equation}
\sum_{j\in S}\sum_{l:(i,l)\in\mathcal{E}}y_{il}^j \geq 
\begin{cases}
0 & \forall i\in\mathcal{V}\setminus\bar{\mathcal{V}} \\
1 & \forall i\in\bar{\mathcal{V}}
\end{cases}
\end{equation}
% \begin{equation}
% \sum_{j\in S}\sum_{l:(i,l)\in\mathcal{E}}y_{il}^j\geq\begin{cases}0&\forall i\in\mathcal{V}\setminus\bar{\mathcal{V}}\\1&\forall i\in\bar{\mathcal{V}}&\end{cases}
% \end{equation} 
对每个单独的传感器节点$j$，进入区域网格节点的次数必须等于离开该节点的次数，以保证路径的连贯性 ：
\begin{equation}
\sum_{l:(i,l)\in\mathcal{E}}y_{il}^j-\sum_{l:(l,i)\in\mathcal{E}}y_{li}^j=0,\quad\forall i\in\mathcal{V},\forall j\in S
\end{equation} 
保证起点约束，即保证传感器节点从其真实位置触发：
\begin{equation}
\sum_{l:(s_k^j,l)\in\mathcal{E}}y_{s_k^jl}^j=1,\quad\forall j\in S
\end{equation} 
为了防止生成与起点不连通的局部环路，施加子环路消除约束：
\begin{equation}
u_i^j-u_l^j+|\mathcal{V}|\cdot y_{il}^j\leq|\mathcal{V}|-1,\quad\forall i\neq l,\forall j\in\mathcal{S}
\end{equation} 
\begin{equation}
1\leq u_i^j\leq|\mathcal{V}|,\quad\forall i\in\mathcal{V},\forall j\in\mathcal{S}
\end{equation} 
变量约束：
\begin{equation}
y_{il}^j\in\{0,1\}
\end{equation} 

该问题本质上属于多旅行商问题,为NP-hard问题。为了降低计算复杂度。
本文采用贪心策略的近似求解方法。

定义$\mathrm{U}$为动态更新的待搜索区域集合，初始时刻$U=\bar{\mathcal{V}}$
将传感器$j$规划路径的尾节点记为$v_{last}^j$，初始时刻，$v_{last}^j$为传感器的位置。
集合$\mathcal{S}$中的传感器依次迭代式选择新增节点，对于任一传感器$j$，其局部决策目标
是从当前未访问集合$\mathrm{U}$中，寻找距离$v_{last}^j$移动代价最小的节点$v^{*}$，
采用欧氏距离作为节点间的移动代价度量。当传感器$j$选定最优节点$v^{*}$后，待搜索区域集合更新过程如下：
首先将$v^{*}$加入传感器$j$的待搜索序列中：
\begin{equation}
\bar{\nu}_j\leftarrow\bar{\nu}_j\cup\{v^*\}
\end{equation} 
更新传感器$j$的路径尾节点，使其指向最新添加的节点：
\begin{equation}
v_{last}^j\leftarrow v^*
\end{equation} 
从待搜索区域集合$\mathrm{U}$中去除该节点，确保其他传感器在后续分配中不再重复计算该节点：
\begin{equation}
U\leftarrow U\setminus\{v^*\}
\end{equation} 
迭代执行上述步骤，直至$U=\emptyset$。
在得到传感器节点的待搜索网格序列后，需要采取控制动作$c_k\in\mathbb{C}_k$，使传感器节点移动到指定路径。
传感器的控制动作目标函数可以表示为
\begin{equation}
c_k^*=\arg\min_{c_k\in\mathbf{C}_k}\xi_{search}(c_k,v)
\end{equation} 
其中，$v\in\bar{\nu}_j$表示该传感器待搜索区域列表列表中下一个待搜索区域的位置，
$\xi_{search}$返回传感器在应用假设的控制动作 $c_k$ 后的预测位置与下一个未访问节点 $v$ 之间的欧氏距离 。

其过程如表4.1所示。

\begin{table}[htbp]
  \centering
  \label{alg:mappo_coop}
  \renewcommand{\arraystretch}{1.2} % 稍微增加行高使得公式显示更美观
  \begin{tabular}{p{0.95\textwidth}}
    \toprule
    \textbf{Algorithm 4.1} 多传感器协同搜索算法 \\
    \midrule
    \textbf{输入：} 搜索传感器集合$\mathcal{S}$以及位置 $s_k^j, j \in \mathcal{S}$，
    区域网格图 $G = (\mathcal{V}, \mathcal{E})$，待搜索区域集合 $U = \mathcal{V}$ \\
    初始化路径 $T_j \leftarrow s_k^j$，$T_j$ 包含智能体 $j \in \bar{S}$ 的搜索路径 \\
    设置指示变量 $\bar{\delta} = 0$，用于确定首个分配节点的目标序列 \\
    \textbf{当} $|\bar{\mathcal{V}}| \neq \emptyset$ \textbf{时执行：} \\
    \quad (a) 更新步进变量 $\bar{\delta} = \bar{\delta} + 1$，确定当前智能体索引 $j = \bar{\delta} \bmod |\bar{S}|$ \\
    \quad (b) 计算最小代价节点：$\min_{l \in \bar{\mathcal{V}}} c(i, l)$，其中 $i$ 为 $T_j$ 中最后添加的节点 \\
    \quad (c) 更新路径 $T_j \leftarrow l$，并更新待访问集合 $\bar{\mathcal{V}} = \bar{\mathcal{V}} \setminus \{l\}$ \\
    \bottomrule
  \end{tabular}
\end{table}

\section{仿真实验与结果分析}
仿真实验一

仿真实验一的目的是验证任务分解策略在解决搜索与跟踪资源竞争问题上的有效性。

仿真实验一中设置监视区域为$[-2000m,2000m]x[-2000m,2000m]$的正方形区域，
目标的运动模型如3.6小节所述，下面给出

对比算法：算法A（本文方法）：基于任务分配的两阶段多传感器管控算法。
算法B（基线方法）：第三章提出的加权联合优化管控算法（选取最佳权重参数）。
算法C（常规方法）：纯信息驱动管控算法或随机搜索/就近跟踪算法（作为底线参考）。

场景设置：在固定区域内，设置多个目标的动态新生与消亡。

数据展示：绘制三种算法随时间变化的 OSPA距离曲线（反映整体跟踪和势估计性能）。
展示传感器是如何根据环境变化和跟踪质量进行角色切换的。

绘制 累积发现目标数量 或 平均目标发现延迟时间 的柱状图。（目标跟踪时间与目标生命周期的平均比率）

结果分析重点：强调本文方法避免了帕累托最优下的折中妥协，实现了搜索和跟踪子任务的“双赢”。

目标的初始状态及参数如表3.3所示：
\begin{table}[H]
    \centering
    \caption{仿真场景一多目标参数}
    \label{tab:scenario_1}
    % 定义四个居中对齐的列
    \begin{tabular}{cccc}
        \toprule[1.5pt] % 顶部粗线，可按需调整磅值
        仿真目标 & 初始位置（m） & 初始速度（m/s） & 出生/消失时刻（s） \\
        \midrule[0.5pt] % 中间细线
        目标 1 & $[-1500, 1000]$ & $[15, 0]$ & 20/100 \\
        目标 2 & $[-1100, -800]$ & $[20, 0]$ & 1/90 \\
        目标 3 & $[-1100, -1100]$ & $[20, 0]$ & 1/80 \\
        目标 4 & $[-1100, -900]$ & $[20, 0]$ & 1/80 \\
        \bottomrule[1.5pt] % 底部粗线
    \end{tabular}
\end{table}

两种方法下的OSPA值图
所提方法在不同时刻的传感器状态
两种方法的计算效率



实验二：算法计算效率与可扩展性分析实验目的：
验证本文方法在降低计算复杂度、提升系统扩展性方面的优势。

场景设置：保持监视区域和目标分布规律不变，
逐步增加传感器网络中的节点规模（例如 $N = 4, 8, 12, 16, 20$）。

数据展示：绘制单步决策平均耗时随传感器数量变化的曲线图（横轴为传感器数量，纵轴为对数坐标系下的计算时间）。
总仿真中对整个区域的搜索率

结果分析重点：指出第三章方法因联合动作空间导致“维度灾难”，
而本章的分组优化机制极大地降低了计算负担，满足大规模传感器网络的实时管控需求。
